\subsection{Good unpointed twists and the prescribed-point obstruction}
\label{sec:twists-pointing-obstruction}

\paragraph{Finding A48 (the exact limit of the constant-twist supply argument).}
Let $C/k$ be a smooth projective geometrically connected curve of
genus $g\ge2$ over a finite field of characteristic $p$.
Let $D\to C$ be a geometrically connected finite etale cyclic
cover of odd degree $n>1$ prime to $p$, with constant deck group
$G=\langle\sigma\rangle$ over $k$. Let $P_n$ be its actual
primitive conductor-$n$ abelian subvariety, as in A46, and put
\[
                       f_n=\dim_{\mathbf F_p}P_n[p](\bar k).
\]
Among the $n$ constant twists $D_t\to C$, indexed by
$t\in\mathbf Z/n$, at least
\[
                n-f_n\ \ge\ n-(g-1)\varphi(n)
 \tag{TW.1}
\]
have primitive factor with no nonzero rational $p$-power
torsion. In genus two the lower bound is at least
$n-\varphi(n)>0$, for every such degree $n$.
The conclusion concerns \emph{unpointed} covers.

For a specified $P\in C(k)$, exactly one of these twists is
$P$-split. If that pointed twist has nonzero rational primitive
$p$-primary torsion, then every other twist has nonzero rational
primitive $p$-primary torsion over every constant extension in
which its fibre above $P$ acquires a rational point.
The same no-improvement assertion holds for twists of a fixed
geometrically cyclic cover with nonconstant deck group, starting
from its pointed model over $k$.
Thus good unpointed twists, followed by a pointing extension,
do not by themselves supply the witness required by A46.

\begin{proof}
\emph{1. Bound the primitive character multiplicities.}
Over $\bar k$, the coherent eigensheaves of the cyclic cover are
nontrivial degree-zero line bundles for every nontrivial
character. Riemann--Roch gives
\[
             \dim_{\bar k}H^1(C_{\bar k},L_\chi)=g-1.
 \tag{TW.2}
\]
The functorial Artin--Schreier stable-part comparison embeds
$H^1_{\mathrm{et}}(D_{\bar k},\mathbf F_p)\otimes\bar k$
as the Frobenius-stable part of $H^1(D_{\bar k},\mathcal O_D)$.
We use the same precise comparison as A32: Katz,
Proposition~2.0, Corollary~2.1 and Propositions~2.2.4--2.2.5
\cite{KatzCongruence}. These statements apply to proper schemes
in every cohomological degree, with no restriction on genus or
ordinarity. The comparison is compatible with deck automorphisms.
The functorial identification of $H^1_{\mathrm{et}}(D_{\bar k},
\mathbf F_p)$ with the dual of $T_pJ_D/pT_pJ_D$, and
$T_pJ_D/pT_pJ_D\simeq J_D[p](\bar k)$, shows that passing to
the $p$-torsion representation only inverts characters. Therefore,
in
\[
 V=P_n[p](\bar k),\qquad
 V\otimes_{\mathbf F_p}\overline{\mathbf F}_p
                    =\bigoplus_{\chi\ {\rm faithful}}V_\chi,
\]
each multiplicity $m_\chi=\dim V_\chi$ is at most $g-1$.
In particular $f_n=\sum_\chi m_\chi\le(g-1)\varphi(n)$.
The sum is a direct sum because $p\nmid n$; no ordinary-rank
hypothesis has been imposed.

\emph{2. Count the twists with a primitive fixed vector.}
Let $F$ denote arithmetic Frobenius acting on geometric points.
It commutes with $G$, since the deck transformations are defined
over $k$. A constant twist $D_t$ has Frobenius descent action
\[
                           F_t=\sigma^t F.
 \tag{TW.3}
\]
These $n$ twists exist by finite Galois descent: the constant
cocycle sends Frobenius to $\sigma^t$ and is trivial after a
finite constant extension. They are distinct as forms over $C$,
whose full geometric automorphism group over $C$ is $G$.
Their primitive abelian subvarieties have the same geometric
$p$-torsion space $V$ and the twisted action (TW.3).

On $V_\chi$, let $\alpha$ run through the eigenvalues of $F$,
counted with algebraic multiplicity. The action of $F_t$ has
an eigenvalue one on this summand exactly when
\[
                         \chi(\sigma)^t\alpha=1
 \tag{TW.4}
\]
for at least one such $\alpha$. Since $\chi$ is faithful,
$t\mapsto\chi(\sigma)^t$ is injective. Each $\alpha$ therefore
forbids at most one value of $t$. This remains true when the
matrix of $F$ has a nontrivial Jordan block: a fixed vector
exists exactly when one is an eigenvalue. Summing this union
bound over the faithful characters gives at most $f_n$ bad
twists. A finite rational $p$-primary group is nonzero if and
only if it has an element of order $p$. Thus the complementary
twists have zero rational primitive $p$-primary group, proving
(TW.1).

In particular every geometric cyclic cover of a genus-two
curve, after descending it and making its deck group constant,
has a twist with this unpointed rational vanishing property.
This conclusion uses neither an Artin-prime statement nor GRH.
It makes no assertion about the rational fibre above a
prescribed base point.

\emph{3. The point chooses one twist.}
Choose a geometric point $R$ of the fibre above $P$, and write
$F(R)=\sigma^a R$. Because $F$ commutes with $G$, its action
on this entire fibre is translation by $a$. The action for
$D_t$ is translation by $a+t$. It has a fixed point if and
only if $a+t=0$ modulo $n$, in which case every point of the
fibre is rational. Thus the unique $P$-split twist is
$D_P=D_{-a}$, with
\[
                              F_P=\sigma^{-a}F.
\]
Over $k_m=\mathbf F_{q^m}$, the fibre of $D_t$ above $P$
has a rational point exactly when $m(a+t)=0$ modulo $n$.
In that situation
\[
 F_t^m=\sigma^{m(a+t)}F_P^m=F_P^m.
 \tag{TW.5}
\]
The same geometric cover with these equal descent actions is
the same form over $k_m$. Its actual primitive subvariety is
identified as well. Consequently
\[
             P_{n,P}(k)[p^\infty]
                      \hookrightarrow P_{n,t}(k_m)[p^\infty].
 \tag{TW.6}
\]
Any nonzero torsion in the pointed model therefore returns as
soon as the other twist can be used as a pointed cover.
This proves the stated obstruction.

\emph{4. The same cancellation for nonconstant deck groups.}
Start now with a geometrically cyclic pointed cover $D_P/k$
whose deck group is not required to be constant. Fix its chosen
point $R$ over $P$ as the origin in the geometric fibre, and
write
\[
 F_P\sigma F_P^{-1}=\sigma^a,\qquad a\in(\mathbf Z/n)^\times.
\]
For a twist represented by $F_t=\sigma^t F_P$, define
$N_m(a)=1+a+\cdots+a^{m-1}$ in $\mathbf Z/n$. On the
geometric fibre, the $m$th Frobenius power acts by
\[
                         x\longmapsto a^m x+N_m(a)t.
 \tag{TW.7}
\]
If this fibre has a rational point over $k_m$, there is an $x$
with $N_m(a)t=(1-a^m)x$. Conjugating the full descent action
by the geometric deck translation $\sigma^x$ then gives
\[
 \sigma^{-x}F_t^m\sigma^x
       =\sigma^{N_m(a)t-(1-a^m)x}F_P^m=F_P^m.
 \tag{TW.8}
\]
Thus the two covers are isomorphic over $C_{k_m}$, and their
canonical primitive subvarieties are isomorphic over $k_m$.
The nonzero rational $p$-primary subgroup of the pointed model,
if present, persists under this isomorphism and the constant
extension. This proves the final assertion. Unlike (TW.1),
no count of nonconstant twists is being asserted.
\end{proof}

\paragraph{What remains for a prescribed-point witness.}
For a fixed geometric cyclic cover and a prescribed point,
twisting can relocate a Frobenius eigenvalue one while
simultaneously relocating the rational fibre. Equations
(TW.5)--(TW.8) show why that operation cannot create the
pointed vanishing condition from a model which fails it.
This does not rule out changing the geometric cyclic cover,
its conductor, or the base curve. Nor does it rule out direct
Fourier verification of a successful pointed cover, as in the
retained examples.

The good-twist bound is therefore a positive unpointed supply
statement with an exact pointing obstruction. It does not
close the universal rational primitive-factor existence gate.
In particular, an average nonvanishing assertion, or a positive
density of arithmetic candidate primes under number-field
GRH, cannot be inserted into this argument as proof that the
unique prescribed-point twist is among the good ones.
A further condition controlling that particular twist, or a
construction varying the geometric cover while keeping the
pointing, is still needed.
